\chapter{Retrieval Methodologies}
\label{ch:methodology}

\section{Classical Foundations}
\label{sec:classical}

The agentic components do not replace classical retrieval; they schedule it.
Every sub-query the Planner emits is answered by one of the three mechanisms
below, so their blind spots determine what the agent layer has to compensate
for.

\subsection{Sparse Retrieval: BM25}
\label{sec:bm25}

BM25~\cite{robertson2009probabilistic} scores a document against a query by
summing, over query terms, a saturating function of term frequency weighted by
inverse document frequency:
\[
  \mathrm{score}(q,d) \;=\; \sum_{t \in q} \mathrm{IDF}(t)\cdot
  \frac{f(t,d)\,(k_1+1)}
       {f(t,d) + k_1\!\left(1 - b + b\,\dfrac{|d|}{\mathrm{avgdl}}\right)} .
\]
The values used here, $k_1 = 0.9$ and $b = 0.4$, sit below the classical 1.2
and 0.75 on purpose: the indexed unit is a single lead paragraph averaging four
sentences, where repetition carries little extra evidence and length
differences are incidental. BM25 matches rare surface forms exactly, needs no
training and answers in milliseconds. Its weakness is vocabulary mismatch: a
query and a relevant passage sharing no stemmed term score zero against each
other however obviously they concern the same thing.

\subsection{Dense Retrieval}
\label{sec:dense}

A dense retriever embeds queries and passages independently and ranks by
similarity. This work uses \texttt{bge-small-en-v1.5}~\cite{xiao2023cpack}, a
384-dimensional bi-encoder with L2-normalised vectors in an exact FAISS index,
used asymmetrically as its authors intend. Its strength is paraphrase, exactly
what BM25 cannot see; its weakness is the mirror image. Query and passage never
meet, and each is compressed to 384 dimensions before comparison, a bottleneck
that preserves topic far better than identity, so two similarly named people
are liable to be conflated. The two failure modes are close to complementary,
which is the standing argument for combining them.

\subsection{Hybrid Fusion and Reranking}
\label{sec:hybrid}

The two lists are combined by reciprocal rank fusion~\cite{cormack2009reciprocal},
which discards scores and combines positions:
\[
  \mathrm{RRF}(d) \;=\; \sum_{r \in R} \frac{1}{k + \mathrm{rank}_r(d)},
  \qquad k = 60 .
\]
Rank fusion is preferred over score normalisation because BM25 scores are
unbounded and corpus-dependent while cosine similarities occupy a narrow band;
the two share no scale, and a manufactured one makes the fused score depend on
the shape of the candidate list. Fusion improves recall and leaves head
precision to chance, which is what reranking addresses. A cross-encoder,
\texttt{ms-marco-MiniLM-L-6-v2}, runs full attention over each query and
candidate pair; nothing about the passage can be precomputed, so scoring $n$
candidates costs $n$ forward passes. That fixes the depths at 50 before
reranking and 10 after. When the fused top result already leads the runner-up
by a relative margin of 0.15 the fifty passes cannot change the head and are
skipped, a saving counted in the trace rather than assumed.

\section{Why Single-Shot Retrieval Fails on Multi-Hop Needs}
\label{sec:singleshot}

Everything above computes a scoring function $s(q, d)$. The three methods
differ enormously in expressive power but not in their arguments: all assume
that the relevance of a document is determinable from the query and that
document. For a bridge question the assumption is false, and the consequence is
not that these systems perform poorly but that they solve the wrong problem.

Consider question \texttt{5a7793275542992a6e59df04} from the evaluation
sample: \emph{Which country is the firm that owns Babycham located?} The gold
answer is Australia, and the two gold sentences open the passages
\textbf{Babycham} (``\ldots the name is now owned by Accolade Wines'') and
\textbf{Accolade Wines} (``\ldots headquarters in South Australia''). The
first hop flatters BM25: Babycham is a near-unique term in a 66{,}581-passage
corpus, and lexical matching alone puts the passage at rank one. The second hop
is not merely harder. Of the query terms, not one occurs in the second gold
sentence. BM25 has nothing to match. A bi-encoder produces a non-zero
similarity, but the query embedding is dominated by Babycham and a possessive
relation while the passage says wine business, global, South Australia, and
nothing distinguishes it from the corpus's other company profiles.

The reranker does not close the gap, for a reason that survives any pool width
and any model size. Place the Accolade Wines passage in the pool by fiat and
give the cross-encoder unlimited capacity. The information that would justify
a high score, that Accolade Wines owns Babycham, is present in neither element
of the pair it is shown; it is asserted in a third document. The relevance of
the second-hop passage is a conditional quantity
$\mathrm{rel}(q, d_2 \mid d_1)$ depending on a document the ranker was never
given, and more capacity changes the function class, not the arguments. The
query that retrieves Accolade Wines trivially is ``Where is Accolade Wines
located?'', and it cannot be written in advance because that string exists
only as the output of the first retrieval. The missing capability is a control
structure, not a scoring function.

This is not an edge case. Bridge questions are 5{,}918 of HotpotQA's 7{,}405,
and 2Wiki's compositional, inference and bridge-comparison types, 9{,}536 of
12{,}576 questions, share the conditional structure. Comparison questions fail
differently: both entities are named in the question, so both passages are
individually retrievable, and the difficulty is coverage at a fixed cut-off
plus the fact that the answer is an operation over two passages rather than
the content of either. Section~\ref{sec:decomposition} gives the two shapes
different plan topologies accordingly.

\section{Agentic Retrieval}
\label{sec:agentic}

Figure~\ref{fig:architecture} shows the system this section builds up to.
The remedy is a retrieval process that carries state between rounds, which is
the thread through the agentic IR literature: a model that plans, invokes tools
and revises~\cite{zhang2024aiagent,wu2025agentic}, whether over scientific
literature~\cite{lala2023paperqa}, multimodal
collections~\cite{schneider2025collex}, telecommunications
networks~\cite{zhang2025toward} or geospatial
workflows~\cite{lee2025multiagent}. The four agents below cover the five tasks
of the brief's Table~1.

\begin{figure}[!t]
  \centering
  \resizebox{0.78\textwidth}{!}{%
  \begin{tikzpicture}[
    font=\small,
    box/.style={draw, rounded corners=2pt, minimum height=9mm,
                minimum width=18mm, align=center, inner sep=3pt},
    store/.style={draw, dashed, rounded corners=2pt, minimum height=11mm,
                  align=center, inner sep=4pt},
    flow/.style={-{Stealth[length=2mm]}, thick},
    back/.style={-{Stealth[length=2mm]}, very thick, dashed}
  ]
    \node (q) {Question};
    \node[box, right=8mm of q]     (plan) {Planner};
    \node[box, right=20mm of plan] (retr) {Retrieval\\Agent};
    \node[box, right=20mm of retr] (kg)   {KG\\Navigator};
    \node[box, right=20mm of kg]   (syn)  {Synthesizer};
    \node[box, right=20mm of syn]  (ver)  {Verifier};
    \node[right=15mm of ver]       (ans)  {Answer};

    \node[store, below=11mm of retr] (idx)   {BM25 $+$ dense\\$+$ rerank};
    \node[store, below=11mm of kg]   (graph) {Entity graph};

    \draw[flow] (q) -- (plan);
    \draw[flow] (plan) -- node[above, font=\scriptsize, inner sep=2pt] {plan} (retr);
    \draw[flow] (retr) -- node[above, font=\scriptsize, inner sep=2pt] {passages} (kg);
    \draw[flow] (kg)   -- node[above, font=\scriptsize, inner sep=2pt] {evidence} (syn);
    \draw[flow] (syn)  -- node[above, font=\scriptsize, inner sep=2pt] {answer} (ver);
    \draw[flow] (ver)  -- node[above, font=\scriptsize, inner sep=2pt] {accept} (ans);
    \draw[flow] (retr) -- (idx);
    \draw[flow] (kg)   -- (graph);

    \draw[back] (ver.north) .. controls +(0,18mm) and +(0,18mm) ..
      node[above, font=\scriptsize]
      {\textbf{ReplanDirective}, confidence $<0.55$} (plan.north);
  \end{tikzpicture}}
  \caption{The four-agent architecture. Solid arrows are the forward retrieval
  path; dashed boxes are offline artefacts. The dashed backward edge is the
  Verifier-to-Planner feedback loop, the only information that flows against
  the pipeline.}
  \label{fig:architecture}
\end{figure}

\subsection{Query Decomposition and Refinement}
\label{sec:decomposition}

The Planner converts a question into a directed acyclic graph of at most five
sub-queries. Each node carries its text, a dependency list, an intent and an
expected answer type, and may contain placeholders such as
\texttt{\{\{q1.answer\}\}}. Execution proceeds in topological order, and a
node's placeholders are substituted from its dependencies immediately before it
is routed, so the second-hop query is a genuinely different string at
retrieval time from the one the Planner emitted. This is the territory of
Self-Ask~\cite{press2023selfask} and IRCoT~\cite{trivedi2023interleaving},
with four differences. The plan is a validated data structure rather than a
decoding trajectory, and every mechanical repair is recorded. It is a graph,
so a comparison question yields two independent roots retrieved together.
Query expansion rides inside the same planner response, so it costs no
additional call. And the plan can be revised on evidence from a downstream
agent, which no single-pass trajectory can do.

Two decisions followed from measurement. The plan's \texttt{strategy} label is
derived from the shape of the validated graph and overwrites what the model
reported: on eight representative questions the model produced structurally
correct decompositions 8 times out of 8 and labelled the strategy
\texttt{single\_hop} in 7 of them. It builds the graph well and names it badly.
Separately, the model reliably emits a final combiner node restating the
question and depending on every other node; detecting and skipping it removes
a third of all retrieval on a three-node comparison plan.

The Planner degrades along a deterministic ladder rather than failing: a
comparison template, then a two-node bridge template, then an identity plan
routed to hybrid search with reranking. That floor is not the
\texttt{hybrid\_rerank} baseline, and the ablation grid of
Chapter~\ref{ch:implementation} shows why the distinction matters: an identity
plan still passes through synthesis and verification, so the system with the
Planner removed is a one-node agentic system that answers and cites, not a
ranker.

\subsection{Tool-Augmented Retrieval}
\label{sec:tools}

Seven tools are declared in a registry, four over the indexes and three over
the entity graph, and the description that documents each to a human is the
string injected into the selection prompt, so documentation and registry
cannot drift apart.

Selection among the retrieval tools is a seven-rule decision table over
features of the resolved sub-query: token count, quoted spans, capitalised
runs, out-of-vocabulary rate against the BM25 lexicon, minimum document
frequency among content terms, and whether the node has dependencies. A
planner hint is honoured first (R1); an out-of-vocabulary rate above 0.34
routes to dense search (R2); short quoted or proper-noun-dense queries (R3) and
queries containing a term of document frequency three or fewer (R4) route to
BM25, the Babycham case; dependent nodes (R5) and long or entity-free queries
(R6) route to hybrid; and a default row (R7) makes the table total.

One fact about the table changes what the argument is about. In the headline
HotpotQA run the Planner attached a hint to every one of the 854 sub-queries it
emitted, so R1 fired on all 575 routing decisions and R2 through R7 never
executed. All seven rules are implemented and tested; on this data six are
dead code. The planning call the system pays for anyway already carries a
per-node tool choice, and the table's contribution was to avoid asking the
model twice. The saving is real, but the claim that lexical rules decide as a
model would have decided is untested. Reinforcement learning would enter here
and does not, because the only reward derives from gold supporting facts that
exist solely in the evaluation data.

\subsection{Knowledge Graph Traversal}
\label{sec:kg-traversal}

The KG Navigator operates on the entity graph of
Section~\ref{sec:preprocessing}. Linking is an alias-trie longest match and
costs no model call. Traversal is bounded at two hops and twenty-five
neighbours, expanded in deterministic order of descending edge weight. The
highest-value operation is bidirectional breadth-first search: given two or
more linked seeds it expands from each and reports the meeting node as the
bridge entity. Applied to the Babycham example, the alias Accolade Wines
occurring in the Babycham passage is an edge by construction, so the bridge
entity the retrieval stack cannot surface by ranking is recovered by a graph
walk in milliseconds. This is the concrete sense in which traversal is not
decorative: it is a second, structurally different mechanism for the one
operation single-shot retrieval provably cannot perform.

Traversal is also the most explainable component, because the path is the
explanation: every edge stores the sentence asserting it, so a traversal
result is scored by the same metrics as a retrieved passage. The limits are
equally clear. The graph is induced from mention co-occurrence, and 23.2\% of
HotpotQA edges and 21.0\% of 2Wiki edges carry no label beyond
\texttt{mentions}, so it is a navigation aid rather than a knowledge base.

\subsection{Self-Reflective Validation}
\label{sec:verification}

The Verifier receives the Synthesizer's answer with the evidence pool and
decides whether it is supported. Cited identifiers that do not exist are
recorded as hallucinated citations and contribute no support. A natural
language inference model, \texttt{cross-encoder/nli-deberta-v3-base}, is run
with each cited sentence as premise and the answer sentence as hypothesis; the
Synthesizer emits that sentence as a separate field precisely so the hypothesis
is a well-formed declarative sentence. Two further signals are computed,
citation grounding and retrieval agreement, and the four combine as
\[
  \mathrm{conf} = 0.45\,\mathrm{nli} + 0.25\,\mathrm{citation}
                + 0.15\,\mathrm{retrieval} + 0.15\,\mathrm{llm},
\]
renormalised when the model term is absent. That term is gated: the
adjudication call is issued only when the confidence computed without it falls
inside $[0.40, 0.70]$, since confident acceptances and rejections cannot have
their verdict changed by it.

Above the threshold of 0.55 the verdict is accept. Below it, with re-plans
remaining, the Verifier emits a directive carrying what could not be
established, which sub-queries returned nothing useful, and the normalised
text of every sub-query already tried. Without that ban list an 8B model told
only that its previous attempt was insufficient reliably re-emits the same
decomposition and burns both re-plans for nothing. Below the threshold with no
re-plans left the verdict is abstain, a label rather than a refusal: the best
candidate is still returned. Answer selection at termination is the argument
maximising confidence over all cycles, ties broken toward the earlier answer,
so a bad re-plan cannot silently overwrite a good first attempt, and the error
taxonomy of Chapter~\ref{ch:implementation} contains explicit false-accept and
false-reject categories so the loop is falsifiable.

\section{The Role of the Language Model}
\label{sec:llm-role}

The language model is the component that generalises, and everything else
exists to keep the number of calls small. Three decisions genuinely require it:
mapping an unseen compound question onto a dependency structure, composing a
short answer from heterogeneous evidence, and adjudicating the narrow band
where the numeric signals disagree. Everything else is deterministic. The
worst case over a question is 15 calls against a cap of 20, with four to seven
typical, and two calls are reserved so synthesis and verification always run.

Serving the model locally shaped the method, not only the deployment. Prompt
length is a latency cost, which is why the re-plan prompt passes retrieved
titles rather than text. And \texttt{qwen3} is a hybrid reasoning model that
emits a \texttt{<think>} block by default: with thinking enabled and a
128-token cap it spent the entire budget on a 503-character reasoning block
and produced no answer at all. At that budget thinking starves the response,
so it is suppressed for the evaluated grid, and Section~\ref{sec:improvement}
reports what happens when it is switched back on for one agent with a budget
large enough to hold it.

\section{Architectural Summary}
\label{sec:arch-summary}

Read left to right, Figure~\ref{fig:architecture} is a pipeline; the dashed
edge, along which the Verifier returns a directive to the Planner below a
confidence of 0.55, is what makes it an agent. Every other component is
standard machinery arranged carefully, and Chapter~\ref{ch:implementation}
isolates that edge's effect by ablation.
